\section{Modelo clásico de Fay--Herriot}

\subsection{Estructura en dos niveles}

El modelo clásico de Fay--Herriot se construye mediante dos ecuaciones.

\subsubsection{Nivel 1: modelo de muestreo}

\[
y_i = \theta_i + e_i,
\qquad i=1,\dots,m,
\]
donde:
\begin{itemize}
    \item $y_i$ es el estimador directo,
    \item $\theta_i$ es el verdadero parámetro del área,
    \item $e_i$ es el error de muestreo.
\end{itemize}

Se asume que
\[
e_i \sim N(0,D_i),
\]
y además que los errores de muestreo son independientes:
\[
e_i \perp e_j \qquad \text{si } i \neq j.
\]

Este primer nivel modela la incertidumbre proveniente del diseño muestral.

\subsubsection{Nivel 2: modelo de enlace}

\[
\theta_i = x_i^{\top}\beta + u_i,
\]
donde:
\begin{itemize}
    \item $x_i \in \mathbb{R}^{p}$ es el vector de covariables del área $i$,
    \item $\beta \in \mathbb{R}^{p}$ es el vector de parámetros fijos,
    \item $u_i$ es un efecto aleatorio de área.
\end{itemize}

Se asume que
\[
u_i \sim N(0,A),
\]
con
\[
u_i \perp u_j \qquad \text{si } i \neq j.
\]

También se supone independencia entre errores de muestreo y efectos aleatorios:
\[
e_i \perp u_j \qquad \forall i,j.
\]

\subsection{Modelo combinado}

Sustituyendo $\theta_i = x_i^{\top}\beta + u_i$ en $y_i = \theta_i + e_i$, obtenemos
\[
y_i = x_i^{\top}\beta + u_i + e_i.
\]

Este es el \textbf{modelo clásico de Fay--Herriot}.

\subsection{Distribución marginal de $y_i$}

Como
\[
u_i \sim N(0,A), 
\qquad
e_i \sim N(0,D_i),
\]
e independientes, entonces
\[
y_i \sim N(x_i^{\top}\beta,\; A + D_i).
\]

Por lo tanto,
\[
E(y_i)=x_i^{\top}\beta,
\qquad
\operatorname{Var}(y_i)=A+D_i.
\]

\section{Forma matricial}

Definamos
\[
y=
\begin{pmatrix}
y_1\\
\vdots\\
y_m
\end{pmatrix},
\qquad
X=
\begin{pmatrix}
x_1^{\top}\\
\vdots\\
x_m^{\top}
\end{pmatrix},
\qquad
u=
\begin{pmatrix}
u_1\\
\vdots\\
u_m
\end{pmatrix},
\qquad
e=
\begin{pmatrix}
e_1\\
\vdots\\
e_m
\end{pmatrix}.
\]

Entonces el modelo se escribe como
\[
y = X\beta + u + e.
\]

Con
\[
u \sim N(0,\; A I_m),
\qquad
e \sim N(0,\; D),
\]
donde
\[
D = \operatorname{diag}(D_1,\dots,D_m).
\]

Así, marginalmente,
\[
y \sim N(X\beta,\; V),
\qquad
V = A I_m + D.
\]
